% !PRINTED 70
\weberpage{70}

\begin{center}
{\large \S.\ 21.}

\vspace{0.5em}
{\large\gesperrt{Principal Theorems on Determinants.}}
\end{center}

From the concept of the determinant the first theorems useful for application are easily obtained.

If in the product [\S~20, (3)]
\begin{equation}
 M'=\pm a_{\alpha_1}^{(1)}a_{\alpha_2}^{(2)}\cdots a_{\alpha_n}^{(n)}
\tag{1}
\end{equation}
we rearrange the factors, its value is not changed. We can therefore also order the factors in such a way that the lower indices occur in their natural order \(1,2,\ldots,n\). The upper indices will then occur in a certain permuted order; thus \(M'\) will receive the form
\begin{equation}
 \pm a_1^{(\beta_1)}a_2^{(\beta_2)}\cdots a_n^{(\beta_n)},
\tag{2}
\end{equation}
where
\[
 (\beta_1,\beta_2,\ldots,\beta_n)=\mathfrak B
\]
is, just as \((\alpha_1,\alpha_2,\ldots,\alpha_n)=\mathfrak A\), an arrangement of the digits \(1,2,\ldots,n\). One obtains the arrangement \(\mathfrak B\) by undoing, in the factors of \(M'\), the transpositions which led to \(\mathfrak A\), beginning with the last, in order to recover the original arrangement in the row of the lower indices. The resulting order of the upper indices is then the arrangement \(\mathfrak B\). It follows that \(\mathfrak B\) belongs to the first or the second kind according as \(\mathfrak A\) belongs to the first or the second kind, since both arise through the same number of transpositions. The totality of the \(\mathfrak B\)'s represents, just like the totality of the \(\mathfrak A\)'s, all permutations of the \(n\) elements, since two different \(\mathfrak A\)'s can never lead to the same \(\mathfrak B\). Thus is proved:

\gesperrt{I. The determinant \(\Delta\) can also be formed by permuting the upper indices in the principal term \(a_1^{(1)}a_2^{(2)}\cdots a_n^{(n)}\) in all possible ways, assigning to every product so formed the positive or negative sign according as the applied permutation belongs to the first or second kind, and then taking the sum of all these products.}
